% !TeX root = ../深度学习与大语言模型技术原理.tex
% 由主文档通过 \input 载入；请编译主文档。
\setcounter{chapter}{22} % 保留原书章号。
\chapter{FlashAttention、分页存储与动态调度}
\label{ch:efficientattention}

\section{在线 softmax 的精确合并}

普通注意力先生成 $n\times n$ 分数，再写入显存、读取并归一化，产生大量中间数据搬运。考虑一个查询的分数 $s_j$ 与值向量 $v_j$，最终输出为 $\sum_je^{s_j}v_j/\sum_je^{s_j}$。对一组已处理位置，维护最大值 $m$、分母 $\ell=\sum_je^{s_j-m}$ 和未归一化分子 $u=\sum_je^{s_j-m}v_j$。

新块的统计量为 $m_b,\ell_b,u_b$，合并时令
\begin{align}
m'&=\max(m,m_b),\\
\ell'&=e^{m-m'}\ell+e^{m_b-m'}\ell_b,\\
u'&=e^{m-m'}u+e^{m_b-m'}u_b.
\label{eq:online-softmax}
\end{align}
输出为 $u'/\ell'$。这些式子只是把旧、新两部分指数统一到同一个最大值下，因此在实数算术中与一次性 softmax 完全一致。第一块可直接初始化统计量；全被屏蔽的块不参与合并，避免无效的无穷差。

\section{FlashAttention 的内存与计算取舍}

FlashAttention 把查询、键和值分块加载到片上存储，利用式~\eqref{eq:online-softmax} 累积结果，避免将完整分数矩阵写入高带宽显存（high bandwidth memory, HBM）\cite{flash}。反向阶段利用保存的统计量重算必要分块，进一步减少中间状态存储。它改变的是执行与数据搬运方式，而非把稠密注意力替换成某个近似稀疏核。

稠密注意力的主要算术复杂度仍为 $O(n^2d)$。称为精确注意力，是指计算同一数学函数；有限精度下，归约顺序不同仍可能导致数值差异。性能取决于头维度、块大小、掩码、数据类型和硬件，不能把某篇论文的加速倍数无条件用于另一工作负载。

FlashAttention-3 针对特定 GPU 架构进一步利用异步执行和低精度计算组织\cite{flash3}。理解这些改进需要区分算术、存储带宽、内核并行和同步等待；换用较新内核之前，仍应验证支持范围与数值误差。

\section{PagedAttention 与写时复制}

PagedAttention 将请求的逻辑 KV 序列映射到固定大小的物理块，无需提前为每条请求分配最大长度的连续空间\cite{paged}。块表描述逻辑顺序，注意力内核依表读取相关键值。这样可以减少大连续分配产生的外部碎片，并使请求随生成逐块增长。

若每块容纳 $b_t$ 个词元，单个非共享请求最后一块的未使用容量少于 $b_t$ 个词元；实际仍有块表、对齐与共享策略开销。块太大增加尾部浪费，块太小增加管理和访问成本。分页并不减少注意力数学运算，也不消除每次读取历史 KV 的代价。

分支生成或共享前缀可以让多个请求引用相同物理块。共享块需要引用计数；某请求要修改共享尾块时，应写时复制（copy-on-write），避免污染其他分支。缓存的释放、取消和异常路径同样需要维护计数，否则可能泄漏显存或错误复用数据。

\section{连续批处理、前缀缓存与分块预填充}

连续批处理（continuous batching）在生成迭代边界移除已结束请求并加入新请求，减少固定批量等待最长回答的空闲。调度器需要同时考虑计算词元预算、KV 容量、等待时间和公平性；只追求总吞吐可能让长请求或新请求长期等待。

前缀缓存（prefix caching）复用相同前缀的已计算 KV。有效复用要求词元、模型权重、适配器、位置与相关配置一致；语义相似的提示不能直接共享缓存。系统可以用哈希块或基数树管理共享前缀，SGLang 的 RadixAttention 是相关实现思路\cite{sglang}。缓存加速预填充，却仍需为每个请求生成新的后缀。

分块预填充（chunked prefill）把长输入分成若干调度块，让其与其他请求的解码交错执行，降低长预填充造成的阻塞。每块依然需要正确访问前文 KV，因此不会改变因果上下文。块大小折中内核效率、首词元延迟和其他请求的词元间延迟，应在真实输入输出长度分布下评估。

